1. Two tracks, your team owns DAQ
DAQ (yours): LLM agents for live runtime systems — Run Control, Process Management, Configuration, Information Service, Error Reporting

2. Your work is genuinely greenfield
No prior MCP framework exists for the five live DAQ services — you're building the first one
3. It's 5 separate agents, not 1 agent for everything
One agent per service (5 total), each built on a shared MCP framework/pattern
One LLM orchestration layer sits on top, routes the shifter's question to the right agent(s), stitches the answer together
From the shifter's side it feels like one assistant; architecturally it's 5 modular agents + 1 orchestrator
4. The four-stage build (WP1→WP4)
WP1: study the 5 services, how they expose data
WP2: design the one reusable MCP pattern all 5 agents will follow (your core contribution)
WP3: build + validate the first agent (one service) to prove the pattern works
WP4: reuse the framework to build the remaining four agents
5. Deliverables
Architecture document (agent-to-service protocols + orchestration layer design)
5 working MCP agents (code)
Per-agent benchmarking report — 5 separate report cards, not one combined evaluation
6. Benchmarking = the part to plan early
Five dimensions, per agent: reliability, robustness, output consistency, performance, resource usage
7. Publication plan
Consider splitting into two papers: (a) architecture/framework across 5 agents (b) Benchmarking LLM agents on live scientific infrastructure
8. Who will use our System : Shifters and Experts
Shifters / Operators run the live system day-to-day, not just during. They need routine help: checking status, following procedures, finding documentation
When something unusual happens, they currently escalate to an expert — that's the slow part this assistant reduces
Simple idea: the assistant is a natural-language front door to the whole system — for routine questions and for when something's wrong


